\documentclass[../DoAn.tex]{subfiles}
\usepackage[utf8]{vietnam}
\usepackage{amsmath}
\usepackage{amsfonts}
\usepackage{geometry}
\usepackage{listings}
\usepackage{xcolor}

\geometry{a4paper, margin=1in}

\lstset{
    basicstyle=\ttfamily\small,
    breaklines=true,
    frame=single,
    numbers=left,
    numberstyle=\tiny,
    keywordstyle=\color{blue},
    stringstyle=\color{red},
    commentstyle=\color{olive}
}

\begin{document}

% Defining the chapter introduction
\section{Mở đầu chương}
Trong bối cảnh công nghệ thông tin phát triển mạnh mẽ và nhu cầu số hóa giáo dục ngày càng tăng cao, việc ứng dụng trí tuệ nhân tạo (AI) vào các hoạt động giảng dạy và hỗ trợ học tập đã trở thành xu thế tất yếu. Một trong những vấn đề phổ biến mà các giảng viên thường gặp phải là việc phải trả lời lặp đi lặp lại các câu hỏi đơn giản, cơ bản từ sinh viên qua email, điều này không chỉ tốn thời gian mà còn làm giảm hiệu quả công việc nghiên cứu và giảng dạy.

Chương này trình bày tổng quan về giải pháp được đề xuất trong đồ án: xây dựng hệ thống chatbot ứng dụng xử lý ngôn ngữ tự nhiên (NLP) để hỗ trợ giảng viên tự động trả lời các câu hỏi đơn giản của sinh viên gửi qua email. Giải pháp này không chỉ giúp tối ưu hóa thời gian làm việc của giảng viên mà còn đảm bảo sinh viên nhận được phản hồi nhanh chóng và chính xác.

Nội dung chương bao gồm việc mô tả kiến trúc tổng thể của hệ thống, các thành phần chính và cách thức hoạt động của từng module, luồng xử lý dữ liệu từ đầu vào đến đầu ra, cũng như các công nghệ và thuật toán được sử dụng. Mục tiêu là giúp người đọc có cái nhìn toàn diện về quy trình hoạt động của hệ thống trước khi đi vào phân tích từng bước cụ thể trong các chương tiếp theo.

% Describing the system architecture and workflow
\section{Kiến trúc tổng thể và luồng hoạt động}

\subsection{Tổng quan về kiến trúc hệ thống}
Hệ thống chatbot được thiết kế theo kiến trúc module hóa, bao gồm ba tầng chính:
\begin{itemize}
    \item \textbf{Tầng thu thập và xử lý dữ liệu (Data Layer):} Chịu trách nhiệm truy xuất, làm sạch và tiền xử lý dữ liệu email.
    \item \textbf{Tầng huấn luyện mô hình (Model Layer):} Thực hiện việc huấn luyện và tối ưu hóa mô hình ngôn ngữ.
    \item \textbf{Tầng ứng dụng (Application Layer):} Cǫng cấp giao diện và xử lý các yêu cầu từ người dùng.
\end{itemize}

\subsection{Luồng hoạt động chính}
Giải pháp được đề xuất bao gồm ba giai đoạn chính, được thực hiện tuần tự:
\begin{enumerate}
    \item \textbf{Thu thập và xử lý dữ liệu từ email:} Truy xuất toàn bộ email học vụ của giảng viên từ Gmail API trong khoảng thời gian xác định, sau đó tiến hành làm sạch nội dung, loại bỏ các thành phần không cần thiết, và trích xuất các cặp câu hỏi - câu trả lời để tạo thành tập dữ liệu huấn luyện chất lượng cao.
    \item \textbf{Huấn luyện mô hình chatbot với dữ liệu đã xử lý:} Sử dụng mô hình TinyLLaMA-1.1B làm backbone, kết hợp với kỹ thuật LoRA (Low-Rank Adaptation) để fine-tune mô hình trên tập dữ liệu đã được chuẩn hóa. Quá trình này được tối ưu hóa để có thể chạy trên máy cá nhân với tài nguyên hạn chế.
    \item \textbf{Triển khai chatbot để sinh phản hồi tự động:} Chatbot sau khi được huấn luyện sẽ được tích hợp vào một pipeline hoàn chỉnh có khả năng tiếp nhận câu hỏi mới từ sinh viên và sinh ra câu trả lời phù hợp dựa trên kiến thức đã học được từ dữ liệu email của giảng viên.
\end{enumerate}

% Outlining technologies and tools
\subsection{Công nghệ và công cụ sử dụng}
Hệ thống được xây dựng dựa trên các công nghệ và thư viện hiện đại:
\begin{itemize}
    \item \textbf{Gmail API:} Để truy xuất dữ liệu email một cách an toàn và hiệu quả.
    \item \textbf{PostgreSQL:} Hệ quản trị cơ sở dữ liệu để lưu trữ và quản lý dữ liệu đã xử lý.
    \item \textbf{TinyLLaMA:} Mô hình ngôn ngữ lớn được tối ưu hóa cho môi trường tài nguyên hạn chế.
    \item \textbf{LoRA:} Kỹ thuật fine-tuning hiệu quả cho mô hình ngôn ngữ lớn.
    \item \textbf{Python:} Ngôn ngữ lập trình chính với các thư viện như \texttt{transformers}, \texttt{torch}, \texttt{pandas}.
\end{itemize}

% Detailing email data extraction and processing
\section{Trích xuất và xử lý dữ liệu từ email}

\subsection{Kết nối và xác thực với Gmail API}
Bước đầu tiên là thiết lập kết nối an toàn với Gmail API, sử dụng OAuth 2.0 để xác thực, đảm bảo quyền truy cập hợp lệ vào hộp thư của giảng viên. Quá trình này bao gồm:
\begin{itemize}
    \item Đăng ký ứng dụng với Google Cloud Console.
    \item Thiết lập credentials và scope quyền truy cập phù hợp.
    \item Xử lý refresh token để duy trì kết nối dài hạn.
\end{itemize}

\subsection{Lọc và thu thập email}
Hệ thống thực hiện việc lọc email theo nhiều tiêu chí:
\begin{itemize}
    \item \textbf{Lọc theo miền email:} Chỉ giữ lại email từ các miền đáng tin cậy như \texttt{@gmail.com}, \texttt{@sis.hust.edu.vn}, \texttt{@soict.hust.edu.vn}.
    \item \textbf{Lọc theo thời gian:} Thu thập email trong khoảng thời gian xác định (thường là 1-2 năm gần nhất).
    \item \textbf{Lọc theo nội dung:} Loại bỏ các email spam, quảng cáo, thông báo tự động và chỉ giữ lại các email có tính chất học thuật, hỏi đáp về môn học.
\end{itemize}

\subsection{Tiền xử lý nội dung email}
Quá trình tiền xử lý nội dung email bao gồm:
\begin{itemize}
    \item \textbf{Làm sạch định dạng HTML:} Loại bỏ các thẻ HTML, CSS inline và JavaScript, chỉ giữ lại văn bản thuần túy.
    \item \textbf{Xử lý chữ ký email:} Tự động phát hiện và loại bỏ chữ ký dựa trên các pattern phổ biến như ``Best regards'', ``Sent from my iPhone''.
    \item \textbf{Loại bỏ nội dung trích dẫn:} Xóa bỏ các phần trích dẫn từ email trước đó (thường bắt đầu bằng ``>'' hoặc ``On ... wrote:'').
    \item \textbf{Chuẩn hóa encoding:} Chuyển đổi tất cả văn bản sang UTF-8 để tránh lỗi hiển thị ký tự đặc biệt.
\end{itemize}

\subsection{Phân loại và tổ chức dữ liệu}
Sau khi làm sạch, hệ thống thực hiện:
\begin{itemize}
    \item \textbf{Tách riêng các chuỗi hội thoại (thread):} Sử dụng thread ID và subject line để nhóm các email liên quan.
    \item \textbf{Sắp xếp theo thời gian:} Các email trong mỗi thread được sắp xếp theo thứ tự thời gian.
    \item \textbf{Phân loại vai trò người gửi:}
    \begin{itemize}
        \item Nếu là sinh viên (domain \texttt{@sis.hust.edu.vn} hoặc format đặc biệt) $\Rightarrow$ lưu vào bảng \texttt{questions}.
        \item Nếu là giảng viên (domain \texttt{@soict.hust.edu.vn} hoặc được định nghĩa trước) $\Rightarrow$ lưu vào bảng \texttt{answers}.
    \end{itemize}
\end{itemize}

\subsection{Lưu trữ dữ liệu}
Toàn bộ dữ liệu được lưu trữ trong PostgreSQL với cấu trúc bảng tối ưu:
\begin{itemize}
    \item \textbf{Bảng emails:} Lưu trữ thông tin cơ bản của email (ID, subject, timestamp, sender, recipient).
    \item \textbf{Bảng questions:} Chứa các câu hỏi đã được xử lý từ sinh viên.
    \item \textbf{Bảng answers:} Chứa các câu trả lời tương ứng từ giảng viên.
    \item \textbf{Bảng threads:} Quản lý mối quan hệ giữa các email trong cùng một cuộc hội thoại.
\end{itemize}

% Explaining TinyLLaMA model training
\section{Huấn luyện mô hình chatbot TinyLLaMA}

\subsection{Chuẩn bị dữ liệu huấn luyện}
Dữ liệu câu hỏi - trả lời được kết hợp thành tập huấn luyện định dạng JSONL:
\begin{lstlisting}[language=json,frame=none]
{
  "instruction": "Câu hỏi của sinh viên",
  "output": "Câu trả lời của giảng viên"
}
\end{lstlisting}
Định dạng này đảm bảo tính đơn giản và tương thích với các framework huấn luyện.

\subsection{Tiền xử lý dữ liệu}
Quá trình tiền xử lý bao gồm:
\begin{itemize}
    \item \textbf{Định dạng lại văn bản:} Chuyển đổi sang định dạng conversation với nhãn \texttt{User} và \texttt{Assistant}.
    \item \textbf{Token hóa và chuẩn hóa độ dài:} Sử dụng tokenizer của TinyLLaMA, cắt ngắn hoặc padding để đảm bảo độ dài tối đa 512 tokens.
    \item \textbf{Augmentation dữ liệu:} Áp dụng paraphrasing, synonym replacement để tăng tính đa dạng.
    \item \textbf{Phân chia tập dữ liệu:} Chia thành tập huấn luyện (80\%), validation (10\%) và test (10\%).
\end{itemize}

\subsection{Cấu hình mô hình TinyLLaMA}
\begin{itemize}
    \item \textbf{Lựa chọn mô hình base:} TinyLLaMA-1.1B-Chat-v1.0, đã được pre-train và instruction-tuned.
    \item \textbf{Cấu hình LoRA:}
    \begin{itemize}
        \item rank (r): 16
        \item alpha: 32
        \item dropout: 0.1
        \item target\_modules: \texttt{["q\_proj", "v\_proj", "k\_proj", "o\_proj"]}
    \end{itemize}
\end{itemize}

\subsection{Quá trình huấn luyện}
\begin{itemize}
    \item \textbf{Hyperparameters:}
    \begin{itemize}
        \item Learning rate: 2e-4 với cosine scheduler.
        \item Batch size: 4 (gradient accumulation: 4).
        \item Epochs: 3-5.
        \item Optimizer: AdamW với weight decay 0.01.
    \end{itemize}
    \item \textbf{Monitoring và evaluation:} Theo dõi loss, perplexity, sử dụng early stopping.
    \item \textbf{Checkpointing:} Lưu checkpoint sau mỗi epoch.
\end{itemize}

\subsection{Tối ưu hóa cho môi trường cá nhân}
\begin{itemize}
    \item Sử dụng mixed precision training (FP16).
    \item Gradient checkpointing để tiết kiệm bộ nhớ.
    \item Batch size nhỏ với gradient accumulation.
\end{itemize}

% Summarizing results and future directions
\section{Kết quả đạt được}

\subsection{Tầng thu thập và xử lý dữ liệu (Data Layer)}
\begin{itemize}
    \item \textbf{Module kết nối Gmail API:} Xác thực OAuth 2.0 và truy xuất email thành công.
    \item \textbf{Module xử lý dữ liệu email:} Làm sạch HTML, chữ ký, nội dung trích dẫn, chuẩn hóa encoding.
    \item \textbf{Hệ thống lưu trữ PostgreSQL:} Thiết lập bảng tối ưu để quản lý dữ liệu.
\end{itemize}

\subsection{Tầng huấn luyện mô hình (Model Layer)}
\begin{itemize}
    \item \textbf{Chuẩn bị dữ liệu huấn luyện:} Chuyển đổi dữ liệu email thành định dạng JSONL.
    \item \textbf{Cấu hình và huấn luyện TinyLLaMA:} Fine-tune thành công với LoRA trong 3 epoch.
    \item \textbf{Tối ưu hóa cho máy cá nhân:} Áp dụng mixed precision training và gradient accumulation.
\end{itemize}

% Outlining future development
\section{Hướng phát triển tiếp theo}

\subsection{Tầng ứng dụng (Application Layer)}
\begin{itemize}
    \item Xây dựng REST API hoặc gRPC service để xử lý yêu cầu từ client.
    \item Thiết kế giao diện web hoặc mobile app cho sinh viên và giảng viên.
    \item Tích hợp module tự động đọc email mới và sinh phản hồi.
\end{itemize}

\subsection{Đánh giá và tối ưu hóa}
\begin{itemize}
    \item \textbf{Đánh giá chất lượng mô hình:} Sử dụng BLEU, ROUGE để đo lường chất lượng phản hồi.
    \item \textbf{Human evaluation:} Thu thập feedback từ giảng viên và sinh viên.
    \item \textbf{Continuous improvement:} Xây dựng pipeline cập nhật mô hình định kỳ.
\end{itemize}

\subsection{Tối ưu hóa hiệu suất}
\begin{itemize}
    \item \textbf{Model optimization:} Quantization, pruning để giảm kích thước mô hình.
    \item \textbf{Caching strategy:} Lưu trữ câu trả lời phổ biến để giảm thời gian phản hồi.
    \item \textbf{Scalability:} Thiết kế kiến trúc mở rộng cho khối lượng email lớn.
\end{itemize}

% Discussing challenges and limitations
\section{Thách thức và hạn chế}

\subsection{Thách thức kỹ thuật}
\begin{itemize}
    \item \textbf{Chất lượng dữ liệu:} Email có thể chứa noise, formatting inconsistency.
    \item \textbf{Tài nguyên hạn chế:} Huấn luyện trên máy cá nhân đòi hỏi tối ưu hóa memory.
    \item \textbf{Đa dạng ngôn ngữ:} Xử lý tiếng Việt khó khăn do ít dữ liệu pre-train.
\end{itemize}

\subsection{Hạn chế của hệ thống}
\begin{itemize}
    \item \textbf{Scope giới hạn:} Chỉ trả lời câu hỏi trong phạm vi dữ liệu huấn luyện.
    \item \textbf{Context understanding:} Khó xử lý câu hỏi có reasoning phức tạp.
    \item \textbf{Personalization:} Thiếu khả năng cá nhân hóa phản hồi.
\end{itemize}

% Concluding the chapter
\section{Kết luận chương}
Chương này đã trình bày tổng quan về kiến trúc và quy trình hoạt động của hệ thống chatbot hỗ trợ giảng viên. Đến thời điểm hiện tại, dự án đã hoàn thành hai tầng chính: thu thập và xử lý dữ liệu (Data Layer) và huấn luyện mô hình (Model Layer). Các thành tựu bao gồm pipeline xử lý email tự động, fine-tuning mô hình TinyLLaMA trên máy cá nhân, và cấu trúc dữ liệu JSONL hiệu quả. Các thách thức như xử lý dữ liệu phức tạp và tối ưu hóa tài nguyên đã được giải quyết. Trong các chương tiếp theo, từng thành phần sẽ được trình bày chi tiết với mã nguồn và kết quả thực nghiệm.

\end{document}